%قبل از شروع مسئله ترجیحاً چیزی ننویسید. اما اگر خواستید بنویسید، در قالب لیستی (فقرات و شمارش) بنویسید.
%برای اضافه کردن نمره، بعدی از شروع‌مسئله، [x] را اضافه کنید و به جای x  نمره مورد نظر را قرار دهید. برای زیرمسئله‌ها هم به همین صورت است.
%\شروع‌مسئله[10]

\شروع‌مسئله []
\شروع{زیرمسائل}
\زیرمسئله % مسئلۀ ۱آ
می‌دانیم که پشته‌ها FILO و صف‌ها FIFO هستند. حال فرض کنید که دو پشته داریم که پشته‌ی اول پر و پشته‌ی دوم خالی می‌باشد. حال به ترتیب اجزای پشته‌ی اول را خوانده و در پشته‌ی دوم می‌ریزیم، چون پشته‌ها FILO هستند، اجزایی که زودتر وارد پشته‌ی اول شده‌اند دیرتر وارد پشته‌ی دوم می‌شوند. چون پشته‌ی دوم نیز FILO است، اجزایی که دیرتر وارد شده‌اند زودتر خارج می‌شوند، پس هر عضوی که زودتر وارد پشته‌ی اول شده باشد، بعد از خالی کردن پشته‌ی اول به پشته‌ی دوم زودتر خارج می‌شود که مانند یک FIFO می باشد. حال با کمک این ایده به کمک دو پشته یک صف را طراحی می‌کنیم:

هنگام push به پشته‌ی اول push می‌کنیم. هنگام pop نیز از پشته‌ی دوم pop میکنیم، اما اگر پشته‌ی دوم خالی بود، ابتدا تمام اعضای پشته‌ی اول را وارد پشته‌ی دوم می‌کنیم. سپس از پشته‌ی دوم pop میکنیم. عملیات push چون عملیات push به پشته است از 
$O(1)$
است. عملیات pop شامل حداکثر $n$ عمل pop برای خالی کردن پشته‌ی اول $n$ عمل push برای وارد شدن به پشته‌ی دوم و یک عملیات pop از پشته‌ی دوم می‌باشد. پس از
$O(n)$
می‌باشد. (گر چه پس از یادگیری تحلیل سرشکن می‌توانید بررسی کنید که مرتبه‌ی زمانی pop نیز به طور سرشکن از
$O(1)$
می‌باشد چون هر عضو حداکثر دو بار pop می‌شود.)

کد مربوط به این الگوریتم در ادامه آمده‌است.

\begin{latin}
\begin{lstlisting}
Queue:
	def push(element):
		stack1.push(element)
	
	def pop():
		if not stack2.empty():
			return stack2.pop()
		while not stack1.empty():
			stack2.push(stack1.pop())
		return stack2.pop()
\end{lstlisting}
\end{latin}


\زیرمسئله % مسئلۀ ۱ب
در این قسمت از دو هرم برای خواسته‌ی مسئله استفاده می‌کنیم، ایده‌ی ما به این صورت است که نیمه‌ی بیشتر اعضا را در یک هرم کمینه و نیمه‌ی کمتر اعضا را در یک هرم بیشینه نگه‌داری نماییم.   این ساختار این مزیت را به ما می‌د‌هد که همیشه میانه در یکی از دو ریشه‌ی این هرم‌ها می‌باشد و با
$O(1)$
قابل دسترسی می‌باشد. حال مسئله این است که چگونه باید عملیات insert را در
$O(\log(n)$
انجام دهیم به طوری که دو پشته‌ی با اختلاف اندازه‌ی حداکثر برابر یک داشته باشیم.

برای انجام اینکار کافیست insert را بگونه‌ای انجام دهیم که اگر دو پشته خاصیت بالا را داشتند، بعد insert نیز این خاصیت حفظ شود. چون که حالت ابتدایی داده ساختار نیز با دو پشته‌ی خالی که خاصیت بالا را دارند شروع می‌گردد، همیشه این خاصیت را خواهند داشت.

برای insert ابتدا عضو جدید را با ریشه‌ی دو هرم مقایسه میکنیم.اگر بیشتر از ریشه‌ی هرم کمینه باشد، آن را در هرم کمینه، اگر کمتر از ریشه‌ی هرم بیشینه بود آن را در هرم بیشینه و اگر بین این دو مقدار بود، آن را در یکی از هرم‌های با اندازه‌ی کمتر قرار می‌دهیم. توجه کنید که چون نیمه‌ی بیشتر اعداد در هرم کمینه و نیمه‌ی کمتر در هرم بیشینه بودند، ریشه‌ی هرم کمینه بیشتر از ریشه‌ی هرم بیشینه می‌باشد. پس تنها یکی از سه حالت بالا می‌توانست رخ بدهد. همچنین با روش بالا، خاصیت بیشتر بودن هرم کمینه نسبت به بیشینه همچنان حفظ می‌گردد.
برای مقایسه‌ها و دسترسی به ریشه‌های دو هرم و قرار دادن اجزا در هرم مجموعا
$O(\log(n))$
هزینه دارد.

حال پس از وارد شدن عضو جدید به یک از دو هرم، اندازه‌ی دو هرم حداکثر به اندازه‌ی دو واحد تفاوت خواهد داشت.
اگر این تفاوت به اندازه‌ی یک واحد باشد، به منظور خود رسیده‌ایم اما فرض کنید که این تفاوت اندازه دو واحد باشد، هرم با اندازه‌ی بیشتر را در نظر بگیرید، می‌توان عضو ریشه‌ی این هرم را جدا کرده و آن را به هرم دیگر وارد کرد، با انجام این کار چون که ریشه‌های هرم ها از تمام اجزای هرم بزرگتر یا کوچکتر می‌باشند، هنوز هم درخت کمینه بیشتر از درخت بیشینه می‌باشد. همچنین اندازه ی دو هرم نیز برابر می‌شود و انجام این کار نیز مرتبه‌ی زمانی insert را از
$O(\log(n))$ 
تغییر نخواهد داد.


\پایان{زیرمسائل}


%اگر مسئله چند قسمت داشت به این صورت عمل کنید:
%\شروع{زیرمسائل}
%\زیرمسئله[0] % مسئلۀ ۱آ
%\زیرمسئله % مسئلۀ ۱ب
%\زیرمسئله % مسئلۀ ۱پ
%\زیرمسئله % مسئلۀ ۱ت
%\پایان{زیرمسائل}

\solution{
	پاسخ نمونه.
}




%%%%%%%%%%%%%%%%%%%%%%%%%%%%%%%%%%%%%%%%%%%%%%%%%%%%%
% برخی دستورات مفید برای تمرین‌ها. %
%%%%%%%%%%%%%%%%%%%%%%%%%%%%%%%%%%%%%%%%%%%%%%%%%%%%%

% فرمول‌های نمونه:
%$f(x) = \Theta(x)$
%$T(x, y) \leq \log(x) + 2^y + \binom{2n}{n}$
% {\tt code\_function}


% لیست موارد بدون شماره‌گذاری:
%\شروع{فقرات}
%    \فقره مورد اول 
%        \شروع{فقرات}
%            \فقره مورد اول داخلی 
%                \شروع{فقرات}
%                    \فقره اولین مورد خیلی داخلی 
%                    \فقره دومین مورد خیلی داخلی 
%                \پایان{فقرات}
%            \فقره مورد دوم داخلی 
%        \پایان{فقرات}
%    \فقره مورد دوم 
%\پایان{فقرات}

% لیست موارد شماره‌گذاری شده:
%\شروع{شمارش}
%    \فقره یک مورد 
%    \فقره یک مورد دیگر 
%    \فقره باز هم یک مورد دیگر
%\پایان{شمارش}

% برای نوشتن یک فرمول چند خطی که یک شماره بخوره. برای حالت‌های دیگر به سایت زیر مراجعه کنید:
% https://www.overleaf.com/learn/latex/Aligning_equations_with_amsmath
%\begin{align} 
%    \begin{split}
%        (x+y)^3 &= (x+y)^2(x+y) \\
%                &= (x^2+2xy+y^2)(x+y) \\
%                &= (x^3+2x^2y+xy^2) + (x^2y+2xy^2+y^3) \\
%                &= x^3+3x^2y+3xy^2+y^3
%    \end{split}                                 
%\end{align}

% نوشتن ماتریس و آرایه دوبعدی
%\begin{align}
%    A = 
%    \begin{bmatrix}
%        A_{11} & A_{21} \\
%        A_{21} & A_{22}
%    \end{bmatrix}
%\end{align}

% اضافه کردن تصویر با شرح. برای اضافه کردن تصویر بدون شرح، فقط دستور درج‌تصویر را استفاده کنید.
%\شروع{شکل}[h!]
%\تنظیم‌ازوسط
%  \درج‌تصویر[پهنا=0.5\پهنای‌متن]{img.jpg}
%\شرح{شرح تصویر}
%\برچسب{تصویر:۱}
%\پایان{شکل}
%\پایان{زیرمسائل}

%کد زیر را ببینید.
%%\lstinputlisting{source_filename.py}
%\begin{latin}
%\begin{lstlisting}
%# Import the modules
%import sys
%import random
%
%ans = True
%
%while ans:
%question = raw_input("Ask the magic 8 ball a question: (press enter to quit) ")
%
%answers = random.randint(1,8)
%
%if question == "":
%sys.exit()
%
%elif answers == 1:
%print "It is certain"
%
%elif answers == 2:
%print "Outlook good"
%
%elif answers == 3:
%print "You may rely on it"
%
%elif answers == 4:
%print "Ask again later"
%
%elif answers == 5:
%print "Concentrate and ask again"
%
%elif answers == 6:
%print "Reply hazy, try again"
%
%elif answers == 7:
%print "My reply is no"
%
%elif answers == 8:
%print "My sources say no"
%\end{lstlisting}
%\end{latin}